% chapters/ch01-workflow.tex
\chapter{The Design Workflow}

\begin{learningobjectives}
In this chapter you will learn a repeatable way to design relational databases without relying on SQL.
You will be able to move from a domain description to a schema that is justified by rules and validated by questions.
\end{learningobjectives}

\section{What database design is}
Relational database design is often introduced as “creating tables”. That framing is convenient, but it is not accurate.
Tables are a \emph{representation}. Design is the work that happens before representation becomes safe.

When design is weak, the system still runs, but meaning becomes unstable.
Facts get duplicated, updates require manual discipline, and invalid states slip into the data because nothing prevents them.
In practice, most “data problems” in organisations are design problems: the model did not make the rules explicit.

\begin{definitionbox}
A \term{relational design} is a model of a domain that preserves meaning and supports correctness.
It does this by making identity explicit and by enforcing (or supporting enforcement of) the domain rules as constraints.
\end{definitionbox}

The key idea is simple: data is not just stored, it is \emph{constrained}.
A schema is a statement about what is allowed to exist.
If the schema allows the wrong states, the application will eventually contain wrong states.

\section{The workflow as a discipline}
A beginner often asks: “Where do I start?” The correct answer is: you start with the domain, not with the database.
You describe the world you are modelling, then you force that description to become precise.
Only then do you translate it into relational structures.

\begin{principlebox}
Use this workflow throughout the book:
\[
\text{Story} \rightarrow \text{Business rules} \rightarrow \text{Conceptual model} \rightarrow
\text{Relational schema} \rightarrow \text{Validation}
\]
\end{principlebox}

This workflow is not a bureaucratic ritual.
It exists because each step removes a different kind of ambiguity.
A story is expressive but vague; rules are precise but incomplete; a conceptual model gives structure; the relational schema commits to enforceable choices; validation tests whether you actually captured the domain.

\section{Step 1: Write a small, realistic story}
A story is not a requirements document.
It should be short, concrete, and close to what people say in real life.
The goal is to capture the domain in a way that a reader can challenge.

\begin{examplebox}
\textbf{Clinic story (minimal)}\\
A clinic has doctors and patients.
Patients book appointments with doctors.
Each appointment has a start time.
A doctor cannot have two appointments at the same time.
\end{examplebox}

Stories can be short in other domains too.

\begin{examplebox}
\textbf{University story (minimal)}\\
A university offers courses.
Students enroll in courses each semester.
An enrollment records the date and final grade.
A student cannot enroll in the same course twice in the same semester.
\end{examplebox}

Notice what is \emph{not} in the story: edge cases, exceptional policies, performance concerns, and implementation details.
Those belong later.
If you start with them, you will build complexity before you have correctness.

\section{Step 2: Convert the story into business rules}
Stories are valuable because they contain meaning, but they are dangerous because they allow interpretation.
The next step is to write rules that are testable.
A rule must be something that can be true or false for a given database state.
If a statement cannot be tested, it will not be enforceable, and it will not guide design decisions.

\begin{definitionbox}
A \term{business rule} is a testable statement that must always hold for the system to be correct.
A good rule is unambiguous, stable, and written in plain language.
\end{definitionbox}

From the clinic story we can extract a small set of rules.
These are not “final”; they are the first precise version of the domain.

\begin{examplebox}
\textbf{Business rules (clinic)}\\
R1: A patient can have many appointments.\\
R2: A doctor can have many appointments.\\
R3: An appointment is with exactly one patient.\\
R4: An appointment is with exactly one doctor.\\
R5: Every appointment has a start time.\\
R6: A doctor cannot be double-booked at the same start time.
\end{examplebox}

Two points matter here.
First, the words “many” and “exactly one” are doing real work: they imply cardinality and optionality.
Second, rule R6 is not merely descriptive; it is a constraint that should prevent invalid states.

\begin{deliverablebox}
For each domain chapter, aim for 10--20 rules.
At minimum, include: one uniqueness rule, one optionality rule, and one time-related rule.
If your rules do not include these, your model will usually miss critical constraints.
\end{deliverablebox}

\begin{principlebox}
Write rules before drawing structures.
If a rule is not written, it is not enforceable.
\end{principlebox}

\section{Step 3: Build a conceptual model}
Rules are precise, but they are not structured.
A conceptual model gives structure by identifying what exists and how it connects.

At this stage you decide what the entities are, what attributes belong to them, and which relationships exist between entities.
You also decide cardinalities: whether a relationship is one-to-one, one-to-many, or many-to-many.
These decisions are not about database syntax.
They are about meaning.

For the clinic, the conceptual picture is straightforward.
The domain contains doctors, patients, and appointments.
Appointments connect one doctor to one patient at a specific time.
The relationships “patient has appointments” and “doctor has appointments” are both one-to-many.
That conceptual understanding is the foundation for the relational schema.

\section{Step 4: Translate into a relational schema (still no SQL)}
The relational schema is where you commit.
Entities become relations, identity becomes keys, and relationships become references.
If your conceptual model is clear, this translation is not creative; it is disciplined.

This book uses a compact notation:
\[
\begin{aligned}
\rel{Appointment}(&\underline{appointment\_id}, start\_time, \\
&patient\_id \rightarrow \rel{Patient}, doctor\_id \rightarrow \rel{Doctor})
\end{aligned}
\]
The underline indicates the primary key, and the arrow indicates a reference to another relation.

In the clinic, the schema is small:
\[
\rel{Patient}(\underline{patient\_id}, name, birth\_date)
\]
\[
\rel{Doctor}(\underline{doctor\_id}, name, specialty)
\]
\[
\begin{aligned}
\rel{Appointment}(&\underline{appointment\_id}, start\_time, \\
&patient\_id \rightarrow \rel{Patient}, doctor\_id \rightarrow \rel{Doctor})
\end{aligned}
\]

The important work is not the existence of these relations.
It is the constraints implied by the rules.
For example, rules R3 and R4 imply that the references from \rel{Appointment} to \rel{Patient} and \rel{Doctor} are mandatory.
Rule R6 implies a constraint that prevents two appointments from using the same doctor at the same start time.
In a real database, the enforcement mechanism depends on how time is represented (discrete slots or intervals), but the modelling intention is already clear.

\begin{notebox}
Do not add SQL early to “make it feel real.”
SQL can hide weak modelling decisions behind syntax.
If the meaning is unclear now, SQL will not fix it later.
\end{notebox}

\section{Step 5: Validate the design without SQL}
A schema is not correct because it “looks right”.
It is correct because it supports the rules and the questions the domain must answer.
Validation is how you prevent elegant-looking diagrams from becoming wrong systems.

One form of validation is question-driven.
You write the key questions and check whether the schema can represent their answers cleanly.
For the clinic, typical questions include: listing appointments for a patient, listing a doctor’s schedule for a day, and detecting double bookings.
If these questions require awkward structures or duplicated facts, the design is usually wrong.

A second form of validation is traceability.
You map each rule to the schema element that enforces it or supports enforcing it.

\begin{examplebox}
\textbf{Rule-to-schema traceability (clinic)}\\[4pt]
\begin{tabularx}{\textwidth}{@{}lX@{}}
\toprule
Rule & Design element(s) \\
\midrule
R3 & Appointment.patient\_id is a mandatory reference to Patient \\
R4 & Appointment.doctor\_id is a mandatory reference to Doctor \\
R6 & Constraint preventing duplicate (doctor\_id, start\_time) in Appointment \\
\bottomrule
\end{tabularx}
\end{examplebox}

Traceability is not busywork.
It is how you prove that your schema matches your intended meaning.
If a rule has no corresponding design element, it is either missing from the model or being enforced elsewhere.
Both cases should be explicit.

\section{Common mistakes}
The workflow fails in predictable ways.

The first mistake is to skip the story and jump straight to relations.
This produces schemas that look tidy but do not match real rules.

The second mistake is to write vague rules like “a patient has appointments.”
Vague rules do not tell you whether appointments are mandatory, unique, or time-bound.

The third mistake is to treat validation as optional.
If you do not test rules against the schema, you are guessing, not designing.

The fourth mistake is to add technical details too early.
Performance and tooling matter later, but correctness must come first.

\section{Mini checklist}
Use this list before you move to the next chapter:
\begin{itemize}[leftmargin=*,nosep]
  \item Do you have a short story in plain language?
  \item Do you have 10--20 testable business rules?
  \item Are entities and relationships named and consistent?
  \item Do you have a compact schema with keys and references?
  \item Can you trace the most important rules to the schema?
\end{itemize}

\begin{chaptersummary}
Relational design becomes manageable when you treat it as a workflow.
You start with a small story, force it into testable rules, structure it with a conceptual model, and translate it into a schema with keys and references.
Finally, you validate the result by checking that rules are traceable and key questions are representable.
This is how design becomes defensible, even before writing a single SQL statement.
\end{chaptersummary}

\begin{exercisebox}
Answer in full sentences.
\begin{enumerate}[label=\arabic*.]
  \item Write a short story (6--8 sentences) for a domain you know.
  \item Extract 12 testable business rules from the story.
  \item Classify each rule as identity, relationship, domain, or business constraint.
  \item List the entities and relationships in plain language with cardinalities.
  \item Write a first relational schema using the book notation.
  \item Create a traceability map for at least 6 rules.
  \item State two key questions your schema must answer and explain where the data lives.
\end{enumerate}
\end{exercisebox}

